\chapter{Fazit und Ausblick}

Dieses Kapitel fasst die zentralen Ergebnisse zusammen, beantwortet die Forschungsfrage und skizziert Weiterentwicklungsmoeglichkeiten.

\section{Zusammenfassung der Ergebnisse}

Gegenstand dieser Arbeit war die Konzeption und Implementierung einer interaktiven Webanwendung zur Veranschaulichung praeemptiver CPU-Scheduling-Verfahren. Ausgangspunkt war die didaktische Herausforderung, dass zeitliche Ablaufentscheidungen in statischen Darstellungen nur eingeschraenkt nachvollziehbar sind.

Im Ergebnis wurde ein Scheduling-Visualizer entwickelt, der mehrere Verfahren in einer konsistenten Simulationsumgebung ausfuehrt und deren Verlaeufe ueber eine interaktive Zeitachsenvisualisierung vergleichbar macht. Ein wesentlicher Beitrag liegt in der klaren Trennung von Simulationskern und Darstellung: Fachlogik, Ereignisprotokoll, Segmente und Metriken werden aus einem einheitlichen, deterministischen Modell abgeleitet. Dadurch sind sowohl Reproduzierbarkeit als auch Nachvollziehbarkeit sichergestellt.

Zusatzlich wurde die Anwendung entlang didaktischer und HCI-orientierter Anforderungen gestaltet. Damit verbindet die Arbeit fachliche Korrektheit der Verfahren mit einer benutzbaren, analysierbaren und fuer Lernkontexte geeigneten Interaktion.

\section{Beantwortung der Forschungsfrage}

Die leitende Forschungsfrage lautete sinngemaess, ob und wie eine interaktive Visualisierung das Verstaendnis praeemptiver Scheduling-Verfahren verbessern und deren Unterschiede nachvollziehbar machen kann.

Auf Basis der konzeptionellen Herleitung, der konkreten Implementierung und der quantitativen Evaluationslogik kann diese Frage bejaht werden. Die Anwendung macht zeitliche Entscheidungsablaeufe sichtbar, erlaubt kontrollierte Parameteraenderungen und liefert vergleichbare Kennzahlen pro Verfahren. Insbesondere die Kombination aus Verlaufssicht und Metrikvergleich reduziert die Distanz zwischen theoretischer Beschreibung und beobachtbarem Verhalten.

Die Arbeit zeigt damit, dass interaktive Algorithmusvisualisierung im Scheduling-Kontext nicht nur demonstrativ, sondern analytisch einsetzbar ist, sofern Reproduzierbarkeit, klare Datenmodelle und transparente Interaktion systematisch umgesetzt werden.

\section{Ausblick auf Weiterentwicklungen}

Aus den Ergebnissen ergeben sich mehrere naheliegende Weiterentwicklungen:

\begin{itemize}
	\item \textbf{Erweiterung des Verfahrensspektrums:} zusaetzliche Scheduling-Strategien und Varianten, insbesondere fuer spezifische Lastprofile.
	\item \textbf{Vertiefte Evaluationsdaten:} breiteres Szenarienset und feinere Parameterstudien, um Sensitivitaeten systematischer zu quantifizieren.
	\item \textbf{Didaktische Zusatzfunktionen:} gefuehrte Aufgabenmodi, integrierte Erklaertexte und adaptive Hinweislogik fuer Selbstlernphasen.
	\item \textbf{Ausbau der Accessibility:} weitere Verbesserungen bei Tastaturbedienung, Screenreader-Feedback und individualisierbarer Darstellung.
	\item \textbf{Export und Reproduktionsartefakte:} strukturierter Export von Szenarien, Runs und Ergebnistabellen zur Nutzung in Lehre und Forschung.
\end{itemize}

Darueber hinaus kann die Anwendung als Ausgangspunkt fuer vergleichende Studien dienen, in denen Lernwirkung, Nutzungsverhalten und Interpretationsqualitaet zwischen verschiedenen Visualisierungsformen empirisch untersucht werden.
